% Included by build-pdf.sh via --include-in-header. Fixes long inline code spans (function
% names, file paths like `crates/dae-runtime/src/block_graph.rs:2416`, CLI flags) overflowing
% the page margin and rendering as invisible/clipped text -- plain `\texttt{}` has no
% hyphenation points, so xelatex can't break a long unbroken identifier at all and just lets it
% run off the page.
%
% Tried first: the `seqsplit` package (inserts break points throughout `\texttt{}` so long
% spans CAN wrap). Abandoned -- it's a "fragile command" that crashes whenever a heading
% containing inline code gets processed in hyperref's PDF-bookmark / TOC moving-argument
% context (two different crash signatures hit during testing, `\pdfstringdefDisableCommands`
% didn't fully resolve it either). Not worth chasing further for a book with ~50 chapters where
% any of them could trigger it.
%
% What's used instead: shrink inline code to \small. This doesn't let a single pathologically
% long identifier wrap mid-token, but it directly shrinks the width of every `\texttt{}` span,
% which resolves the common case (a moderately long function name/path overflowing a line it's
% embedded in) -- and it's a completely standard, robust LaTeX construct with no
% moving-argument fragility.
\let\oldtexttt\texttt
\renewcommand{\texttt}[1]{\small\oldtexttt{#1}}
\sloppy
\emergencystretch=3em

% Same overflow problem for fenced code blocks (```...``` -- the ASCII crate-dependency diagram,
% multi-line netlist/CLI examples): let long lines wrap instead of running off the page.
\usepackage{fvextra}
\fvset{breaklines=true,breakanywhere=true}
